\chapter{Inferencia causal clásica}
\cite{oxChapterOutcome}

\subsection{Estratificación}
Una forma de ajustar a las covariantes es estratificar la muestra y considerar
el efecto del factor de interés $A$ dentro de cada estrato, donde las covariantes
son fijas.
Esta técnica es problemática cuando el número de covariantes es significativo,
puesto que el número de estratos crece exponencialmente en función del número de covariantes,
y el número de muestras por estrato decrece acordemente.
A este concepto se le llama la \say{maldición de la dimensionalidad}\footnotemark{}.
Aplicamos esta técnica en \autoref{ex:efectoConfusor}.

\begin{definicion}[Modificación de efecto] \label{def:modificacion_de_efecto}
    Cuando la relación entre dos variables $X$ e $Y$ está sujeta a una tercera variable
    $L$, de forma que el efecto de $X$ en $Y$ se ve modificado según el valor de $L$,
    decimos que $L$ modifica el efecto de $X$ en $Y$.
    Esto se diferencia del confounding (\autoref{def:confounder}), el cual muestra su efecto
    debido a su modificación en la propensidad de $X$, pero el efecto de $X$ e $Y$ no es
    modificado.
\end{definicion}
\begin{nota}
    Lo anterior no excluye la posibilidad del caso que una variable sea tanto
    confusora como modificadora de efecto al mismo tiempo.
\end{nota}

\footnotetext{O en inglés, \say{the curse of dimensionality}.}

\subsection{Regresión}
Otra forma de ajustar el efecto de un factor de riesgo o tratamiento
al efecto de unas covariantes $X = (X_1, \ldots, X_n)$ es midiendo el efecto con una regresión.

$$
    Y \sim A + \sum_{i=1}^{n} X_i
$$

Esta técnica está sujeta a las suposiciones habituales de regresión, en particular
que los errores tengan una distribución normal multivariante y homocedasticidad de
las variables regresoras (que tengan varianzas similares).
Puede ser un estimador insesgado sobre $Y$ supuesto que $X$ afecte a la distribución de $A$,
eso es, que no haya \emph{modificación del efecto}, con lo cual es apto para estudios randomizados
controlados, pero no para estudios observacionales.
